\documentclass[10pt,a4paper]{article}
\usepackage[utf8]{inputenc}
\usepackage[T1]{fontenc}
\usepackage[ngerman]{babel}
\usepackage{geometry}
\usepackage{xcolor}
\usepackage{tabularx}
\usepackage{array}
\usepackage{booktabs}
\usepackage{enumitem}
\usepackage{helvet}
\usepackage{tikz}
\usepackage{colortbl}
\usepackage{amssymb}
\usepackage{fancyhdr}
\usepackage{paracol}
\usepackage{titlesec}
\usepackage{footmisc}

\renewcommand{\familydefault}{\sfdefault}
\renewcommand{\thefootnote}{\arabic{footnote}}
\geometry{margin=1.5cm, top=2cm, bottom=1.5cm}
\setlength{\parindent}{0pt}
\setlength{\parskip}{4pt}
\setlength{\headheight}{14pt}

% Colors
\definecolor{spanisch}{HTML}{c41e3a}
\definecolor{lightred}{HTML}{fdf2f4}
\definecolor{darkgray}{HTML}{333333}
\definecolor{lightgray}{HTML}{f5f5f5}
\definecolor{english}{HTML}{1a5276}

% Headers
\pagestyle{fancy}
\fancyhf{}
\fancyhead[L]{\small\textcolor{spanisch}{\textbf{Sequenzplan Oberstufe Spanisch M-V}}}
\fancyhead[R]{\small\textcolor{english}{\textbf{Upper Secondary Spanish Sequence Plan M-V}}}
\fancyfoot[C]{\thepage}

% Section formatting
\titleformat{\section}{\large\bfseries\color{spanisch}}{}{0pt}{}
\titleformat{\subsection}{\normalsize\bfseries\color{darkgray}}{}{0pt}{}
\titlespacing{\section}{0pt}{12pt}{6pt}
\titlespacing{\subsection}{0pt}{8pt}{4pt}

% Bilingual command
\newcommand{\bil}[2]{%
\begin{paracol}{2}
\small #1
\switchcolumn
\small\textcolor{english}{#2}
\end{paracol}
}

% Section bilingual
\newcommand{\secbil}[2]{%
\vspace{8pt}
\noindent\colorbox{spanisch}{\parbox{0.48\textwidth}{\color{white}\bfseries #1}}%
\hfill%
\colorbox{english}{\parbox{0.48\textwidth}{\color{white}\bfseries #2}}
\vspace{4pt}
}

\begin{document}

% Title
\begin{center}
\colorbox{spanisch}{\parbox{0.48\textwidth}{\centering\color{white}\Large\bfseries\vspace{4pt}
Sequenzplan Spanisch\\Oberstufe (Kl. 10-12)\\Mecklenburg-Vorpommern
\vspace{4pt}}}%
\hfill%
\colorbox{english}{\parbox{0.48\textwidth}{\centering\color{white}\Large\bfseries\vspace{4pt}
Spanish Sequence Plan\\Upper Secondary (Gr. 10-12)\\Mecklenburg-Vorpommern
\vspace{4pt}}}
\end{center}

\vspace{2mm}
\begin{center}
\small Version 1.0 | Januar 2026 \hfill \textcolor{english}{Version 1.0 | January 2026}
\end{center}

\vspace{4mm}

%% STRUCTURE OVERVIEW
\secbil{1. Struktur der Oberstufe}{1. Upper Secondary Structure}

\bil{%
\textbf{Einführungsphase (Klasse 10):}\\
Greift Sek-I-Kompetenzen auf, bereitet auf Qualifikationsphase vor. Kompensations- und Orientierungsfunktion.

\textbf{Qualifikationsphase (Klasse 11-12):}\\
Vertiefte Allgemeinbildung, wissenschaftspropädeutische Grundbildung, Abitur-Vorbereitung.
}{%
\textbf{Introductory Phase (Grade 10):}\\
Builds on Sec I competencies, prepares for qualification phase. Compensation and orientation function.

\textbf{Qualification Phase (Grades 11-12):}\\
In-depth general education, pre-academic foundation, Abitur preparation.
}

\vspace{2mm}

\begin{center}
\small
\begin{tabular}{|l|c|c|c|}
\hline
\rowcolor{lightgray}
\textbf{Klasse / Grade} & \textbf{GER / CEFR} & \textbf{GK / Basic} & \textbf{LK / Advanced} \\
\hline
10 (Einführung / Intro) & B1 & 3 WStd / 120h & 5 WStd / 200h \\
\hline
11 (Qualifikation 1) & B1+ & 3 WStd / 120h & 5 WStd / 200h \\
\hline
12 (Qualifikation 2) & B2 & 3 WStd / 120h & 5 WStd / 200h \\
\hline
\end{tabular}
\end{center}

%% 4 THEME FIELDS
\secbil{2. Die 4 obligatorischen Themenfelder}{2. The 4 Mandatory Theme Fields}

\vspace{2mm}

\begin{center}
\small
\begin{tabular}{|c|p{5.5cm}|c|c|}
\hline
\rowcolor{lightgray}
\textbf{Nr} & \textbf{Themenfeld / Theme Field} & \textbf{GK} & \textbf{LK} \\
\hline
1 & El individuo y la sociedad \newline \textcolor{english}{\footnotesize The Individual and Society} & 45h & 75h \\
\hline
2 & La identidad nacional y la diversidad cultural \newline \textcolor{english}{\footnotesize National Identity and Cultural Diversity} & 45h & 75h \\
\hline
3 & Aspectos actuales de la política y la sociedad \newline \textcolor{english}{\footnotesize Current Aspects of Politics and Society} & 30h & 50h \\
\hline
4 & Desafíos modernos \newline \textcolor{english}{\footnotesize Modern Challenges} & 30h & 50h \\
\hline
\rowcolor{lightred}
& \textbf{GESAMT / TOTAL} & \textbf{150h} & \textbf{250h} \\
\hline
\end{tabular}
\end{center}

\vspace{2mm}

\bil{%
\textit{Hinweis:} Die verbleibenden ~20\% der Unterrichtszeit stehen für Wiederholung, Vertiefung und schulinterne Schwerpunkte zur Verfügung.
}{%
\textit{Note:} The remaining ~20\% of teaching time is available for review, consolidation, and school-specific focus areas.
}

%% GRADE 10
\secbil{3. Klasse 10 -- Einführungsphase}{3. Grade 10 -- Introductory Phase}

\bil{%
\textbf{Halbjahr 10.1 (August -- Januar)}

\underline{Themensequenz:}
\begin{itemize}[noitemsep,topsep=2pt]
\item Woche 1-2: Diagnose \& Auffrischung
\item Woche 3-6: La juventud de hoy
\item Woche 7-10: Relaciones y convivencia
\item Woche 11-14: Perspectivas futuras
\item Woche 15-16: Wiederholung + Klausur
\end{itemize}

\underline{Grammatik-Schwerpunkte:}
\begin{itemize}[noitemsep,topsep=2pt]
\item Vergangenheitszeiten (Festigung)
\item Subjuntivo Presente (Vertiefung)
\item Konditionalsätze Typ I und II
\end{itemize}
}{%
\textbf{Semester 10.1 (August -- January)}

\underline{Topic Sequence:}
\begin{itemize}[noitemsep,topsep=2pt]
\item Week 1-2: Diagnosis \& Review
\item Week 3-6: Today's Youth
\item Week 7-10: Relationships and Coexistence
\item Week 11-14: Future Perspectives
\item Week 15-16: Review + Written Exam
\end{itemize}

\underline{Grammar Focus:}
\begin{itemize}[noitemsep,topsep=2pt]
\item Past tenses (consolidation)
\item Present Subjunctive (deepening)
\item Conditional sentences Type I and II
\end{itemize}
}

\vspace{2mm}

\bil{%
\textbf{Halbjahr 10.2 (Februar -- Juli)}

\underline{Themensequenz:}
\begin{itemize}[noitemsep,topsep=2pt]
\item Woche 1-4: España: geografía y regiones
\item Woche 5-8: Tradiciones y costumbres
\item Woche 9-12: Latinoamérica: una introducción
\item Woche 13-15: Medios y comunicación
\item Woche 16: \textbf{Sprechprüfung} (B1)
\end{itemize}

\underline{Grammatik-Schwerpunkte:}
\begin{itemize}[noitemsep,topsep=2pt]
\item Subjuntivo Imperfecto (Einführung)
\item Relativsätze
\item Passivkonstruktionen
\end{itemize}
}{%
\textbf{Semester 10.2 (February -- July)}

\underline{Topic Sequence:}
\begin{itemize}[noitemsep,topsep=2pt]
\item Week 1-4: Spain: Geography and Regions
\item Week 5-8: Traditions and Customs
\item Week 9-12: Latin America: An Introduction
\item Week 13-15: Media and Communication
\item Week 16: \textbf{Oral Exam} (B1)
\end{itemize}

\underline{Grammar Focus:}
\begin{itemize}[noitemsep,topsep=2pt]
\item Imperfect Subjunctive (introduction)
\item Relative clauses
\item Passive constructions
\end{itemize}
}

\newpage

%% GRADE 11
\secbil{4. Klasse 11 -- Qualifikationsphase 1}{4. Grade 11 -- Qualification Phase 1}

\bil{%
\textbf{Halbjahr 11.1: El individuo y la sociedad}

\underline{Themensequenz:}
\begin{itemize}[noitemsep,topsep=2pt]
\item Woche 1-3: Jóvenes y pasiones
\item Woche 4-6: Amor y relaciones
\item Woche 7-9: Perspectivas en la adolescencia
\item Woche 10-12: Facetas de la vida en España
\item Woche 13-14: Klausur
\item Woche 15-16: +LK: Formas de convivencia
\end{itemize}

\underline{Leitmedium:} Kurzfilm \textit{Física II}

\underline{Literatur:}
\begin{itemize}[noitemsep,topsep=2pt]
\item \textit{La mirada} (Juan Madrid) -- GK/LK
\item \textit{El método Grönholm} (Galceran) -- \textbf{nur LK}
\end{itemize}

\underline{Kompetenzschwerpunkte:}
\begin{itemize}[noitemsep,topsep=2pt]
\item Filmanalyse, Figurencharakterisierung
\item Tagebucheintrag / Rezension schreiben
\item Diskussion zu Zukunftsperspektiven
\end{itemize}
}{%
\textbf{Semester 11.1: The Individual and Society}

\underline{Topic Sequence:}
\begin{itemize}[noitemsep,topsep=2pt]
\item Week 1-3: Youth and Passions
\item Week 4-6: Love and Relationships
\item Week 7-9: Adolescent Perspectives
\item Week 10-12: Facets of Life in Spain
\item Week 13-14: Written Exam
\item Week 15-16: +Adv: Forms of Coexistence
\end{itemize}

\underline{Lead Medium:} Short film \textit{Física II}

\underline{Literature:}
\begin{itemize}[noitemsep,topsep=2pt]
\item \textit{La mirada} (Juan Madrid) -- Basic/Adv
\item \textit{El método Grönholm} (Galceran) -- \textbf{Adv only}
\end{itemize}

\underline{Competency Focus:}
\begin{itemize}[noitemsep,topsep=2pt]
\item Film analysis, character portrayal
\item Writing diary entries / reviews
\item Discussion on future perspectives
\end{itemize}
}

\vspace{3mm}

\bil{%
\textbf{Halbjahr 11.2: La identidad nacional}

\underline{Themensequenz:}
\begin{itemize}[noitemsep,topsep=2pt]
\item Woche 1-3: Momentos cruciales: El Trialog
\item Woche 4-6: La Guerra Civil y sus secuelas
\item Woche 7-9: Latinoamérica ayer y hoy
\item Woche 10-12: El multiculturalismo
\item Woche 13-14: \textbf{Sprechprüfung} (B1+/B2)
\item Woche 15-16: +LK: España después de la Transición
\end{itemize}

\underline{Leitmedium:} Fidel Castros UNO-Rede 1979

\underline{Literatur:}
\begin{itemize}[noitemsep,topsep=2pt]
\item \textit{El desertador} (Merino) -- GK/LK
\item \textit{Una fosa poco común} (Jambrina) -- GK/LK
\item \textit{Las venas abiertas de A.L.} (Galeano) -- \textbf{nur LK}
\end{itemize}
}{%
\textbf{Semester 11.2: National Identity}

\underline{Topic Sequence:}
\begin{itemize}[noitemsep,topsep=2pt]
\item Week 1-3: Crucial Moments: The Trialogue
\item Week 4-6: Civil War and its Aftermath
\item Week 7-9: Latin America Yesterday and Today
\item Week 10-12: Multiculturalism
\item Week 13-14: \textbf{Oral Exam} (B1+/B2)
\item Week 15-16: +Adv: Spain after the Transition
\end{itemize}

\underline{Lead Medium:} Fidel Castro's UN Speech 1979

\underline{Literature:}
\begin{itemize}[noitemsep,topsep=2pt]
\item \textit{El desertador} (Merino) -- Basic/Adv
\item \textit{Una fosa poco común} (Jambrina) -- Basic/Adv
\item \textit{Las venas abiertas de A.L.} (Galeano) -- \textbf{Adv only}
\end{itemize}
}

\vspace{4mm}

%% GRADE 12
\secbil{5. Klasse 12 -- Qualifikationsphase 2}{5. Grade 12 -- Qualification Phase 2}

\bil{%
\textbf{Halbjahr 12.1: Aspectos actuales}

\underline{Themensequenz:}
\begin{itemize}[noitemsep,topsep=2pt]
\item Woche 1-3: Movimientos migratorios
\item Woche 4-6: La situación de los inmigrantes
\item Woche 7-9: Desarrollos económicos
\item Woche 10-12: +LK: Responsabilidades sociales
\item Woche 13-14: \textbf{Klausur (abiturähnlich)}
\item Woche 15-16: Wiederholung + Vertiefung
\end{itemize}

\underline{Leitmedium:} \textit{La casa en Mango Street} (Cisneros)

\underline{Literatur:}
\begin{itemize}[noitemsep,topsep=2pt]
\item \textit{Abdel} (Enrique Páez) -- GK/LK
\item \textit{La piel de un indio...} (Ribeyro) -- GK/LK
\item Doku: \textit{México-EE.UU.: El gran cruce}
\end{itemize}
}{%
\textbf{Semester 12.1: Current Aspects}

\underline{Topic Sequence:}
\begin{itemize}[noitemsep,topsep=2pt]
\item Week 1-3: Migration Movements
\item Week 4-6: The Situation of Immigrants
\item Week 7-9: Economic Developments
\item Week 10-12: +Adv: Social Responsibilities
\item Week 13-14: \textbf{Exam (Abitur-style)}
\item Week 15-16: Review + Consolidation
\end{itemize}

\underline{Lead Medium:} \textit{La casa en Mango Street} (Cisneros)

\underline{Literature:}
\begin{itemize}[noitemsep,topsep=2pt]
\item \textit{Abdel} (Enrique Páez) -- Basic/Adv
\item \textit{La piel de un indio...} (Ribeyro) -- Basic/Adv
\item Documentary: \textit{México-EE.UU.: El gran cruce}
\end{itemize}
}

\vspace{3mm}

\begin{paracol}{2}
\small
\textbf{Halbjahr 12.2: Desafíos modernos + Abitur}
\par\vspace{1mm}
\underline{Themensequenz:}
\begin{itemize}[noitemsep,topsep=2pt]
\item Woche 1-3: Tendencias ecológicas
\item Woche 4-6: Desafíos de la era digital
\item Woche 7-8: +LK: Sueños y realidades
\item Woche 9-10: Sprechprüfung / Paarprüfung Abitur
\item Woche 11-14: \textbf{Abitur-Vorbereitung}
\end{itemize}
\underline{Leitmedium:} \textit{El agua y la humanidad} (Menchú)
\par\vspace{1mm}
\underline{Literatur:}
\begin{itemize}[noitemsep,topsep=2pt]
\item \textit{547 amigos} (Lorenzo Silva) -- GK/LK
\item \textit{El guardagujas} (Arreola) -- \textbf{nur LK}
\item Karikaturen zu Umweltthemen
\end{itemize}
\switchcolumn
\small\color{english}
\textbf{Semester 12.2: Modern Challenges + Abitur}
\par\vspace{1mm}
\underline{Topic Sequence:}
\begin{itemize}[noitemsep,topsep=2pt]
\item Week 1-3: Ecological Trends
\item Week 4-6: Challenges of the Digital Era
\item Week 7-8: +Adv: Dreams and Realities
\item Week 9-10: Oral Exam / Abitur Pair Exam
\item Week 11-14: \textbf{Abitur Preparation}
\end{itemize}
\underline{Lead Medium:} \textit{El agua y la humanidad} (Menchú)
\par\vspace{1mm}
\underline{Literature:}
\begin{itemize}[noitemsep,topsep=2pt]
\item \textit{547 amigos} (Lorenzo Silva) -- Basic/Adv
\item \textit{El guardagujas} (Arreola) -- \textbf{Adv only}
\item Cartoons on environmental topics
\end{itemize}
\end{paracol}

\newpage

%% EXAMS
\secbil{6. Prüfungsübersicht}{6. Examination Overview}

\vspace{2mm}

\begin{center}
\small
\begin{tabular}{|l|l|l|l|}
\hline
\rowcolor{lightgray}
\textbf{Halbjahr} & \textbf{Prüfung / Exam} & \textbf{Format} & \textbf{Niveau} \\
\hline
10.1 & Klausur 1 / Exam 1 & Lesen + Schreiben & B1 \\
\hline
10.2 & Sprechprüfung 1 / Oral 1 & Tandem, 20 Min. & B1 \\
\hline
11.1 & Klausur 2 / Exam 2 & Lesen + Schreiben + Analyse & B1+ \\
\hline
11.2 & Sprechprüfung 2 / Oral 2 & Tandem, 20 Min. & B1+/B2 \\
\hline
12.1 & Klausur 3 (abiturähnlich) & Alle Kompetenzen & B2 \\
\hline
12.2 & Abitur & Schriftlich + optional mündlich & B2 \\
\hline
\end{tabular}
\end{center}

\vspace{3mm}

\bil{%
\textbf{Abiturprüfung:}
\begin{itemize}[noitemsep,topsep=2pt]
\item \textbf{Schriftlich GK:} 180 Min. (+ 30 Min. Auswahl)
\item \textbf{Schriftlich LK:} 270 Min. (+ 30 Min. Auswahl)
\item \textbf{Mündlich (Option):} Paarprüfung, max. 30 Min.
\item \textbf{AFB-Verteilung:} 20\% / 45\% / 35\%
\end{itemize}
}{%
\textbf{Abitur Examination:}
\begin{itemize}[noitemsep,topsep=2pt]
\item \textbf{Written Basic:} 180 min (+ 30 min selection)
\item \textbf{Written Advanced:} 270 min (+ 30 min selection)
\item \textbf{Oral (Option):} Pair exam, max. 30 min
\item \textbf{AFB Distribution:} 20\% / 45\% / 35\%
\end{itemize}
}

%% GRAMMAR
\secbil{7. Grammatik-Progression}{7. Grammar Progression}

\vspace{2mm}

\begin{paracol}{2}
\small
\textbf{Klasse 10 -- Festigung + Ausbau:}
\begin{itemize}[noitemsep,topsep=2pt]
\item Vergangenheitszeiten (Konsolidierung)
\item Subjuntivo Presente (Vertiefung)
\item Subjuntivo Imperfecto (Einführung)
\item Condicional Simple + Compuesto
\item Oraciones condicionales (Typ I, II, III)
\item Voz pasiva (perifrástica y refleja)
\end{itemize}

\textbf{Klasse 11 -- Vertiefung:}
\begin{itemize}[noitemsep,topsep=2pt]
\item Subjuntivo Imperfecto (Vertiefung)
\item Subjuntivo Pluscuamperfecto
\item Estilo indirecto (vollständig)
\item Conectores textuales (argumentativ)
\item Perífrasis verbales
\end{itemize}

\textbf{Klasse 12 -- Perfektionierung:}
\begin{itemize}[noitemsep,topsep=2pt]
\item Stilistisch differenzierte Anwendung
\item Register (formal/informell)
\item Varietäten (Español americano)
\item Idiomatische Ausdrücke
\end{itemize}

\switchcolumn

\small\color{english}
\textbf{Grade 10 -- Consolidation + Expansion:}
\begin{itemize}[noitemsep,topsep=2pt]
\item Past tenses (consolidation)
\item Present Subjunctive (deepening)
\item Imperfect Subjunctive (introduction)
\item Simple + Compound Conditional
\item Conditional sentences (Type I, II, III)
\item Passive voice (periphrastic and reflexive)
\end{itemize}

\textbf{Grade 11 -- Deepening:}
\begin{itemize}[noitemsep,topsep=2pt]
\item Imperfect Subjunctive (deepening)
\item Pluperfect Subjunctive
\item Reported speech (complete)
\item Textual connectors (argumentative)
\item Verbal periphrases
\end{itemize}

\textbf{Grade 12 -- Perfection:}
\begin{itemize}[noitemsep,topsep=2pt]
\item Stylistically differentiated use
\item Register (formal/informal)
\item Varieties (Latin American Spanish)
\item Idiomatic expressions
\end{itemize}
\end{paracol}

%% LITERATURE MATRIX
\secbil{8. Literatur-Matrix}{8. Literature Matrix}

\vspace{2mm}

\begin{center}
\scriptsize
\begin{tabular}{|p{3.8cm}|p{2.8cm}|p{2cm}|c|c|}
\hline
\rowcolor{lightgray}
\textbf{Titel / Title} & \textbf{Autor / Author} & \textbf{Format} & \textbf{GK} & \textbf{LK} \\
\hline
\multicolumn{5}{|l|}{\cellcolor{lightred}\textbf{Thema 1: El individuo y la sociedad}} \\
\hline
La mirada & Juan Madrid & Kurzgeschichte & $\checkmark$ & $\checkmark$ \\
Desconexión total & Lourdes Miquel & Roman & $\checkmark$ & -- \\
La Casa & Paco Roca & Novela Gráfica & $\checkmark$ & $\checkmark$ \\
El método Grönholm & Jordi Galceran & Theaterstück & -- & $\checkmark$ \\
\hline
\multicolumn{5}{|l|}{\cellcolor{lightred}\textbf{Thema 2: La identidad nacional}} \\
\hline
El desertador & José María Merino & Kurzgeschichte & $\checkmark$ & $\checkmark$ \\
Una fosa poco común & Luis García Jambrina & Kurzgeschichte & $\checkmark$ & $\checkmark$ \\
Las venas abiertas de A.L. & Eduardo Galeano & Sachbuch (Kap.) & -- & $\checkmark$ \\
\hline
\multicolumn{5}{|l|}{\cellcolor{lightred}\textbf{Thema 3: Aspectos actuales}} \\
\hline
Abdel & Enrique Páez & Roman & $\checkmark$ & $\checkmark$ \\
La casa en Mango Street & Sandra Cisneros & Roman & $\checkmark$ & $\checkmark$ \\
La piel de un indio... & J.R. Ribeyro & Kurzgeschichte & $\checkmark$ & $\checkmark$ \\
\hline
\multicolumn{5}{|l|}{\cellcolor{lightred}\textbf{Thema 4: Desafíos modernos}} \\
\hline
547 amigos & Lorenzo Silva & Kurzgeschichte & $\checkmark$ & $\checkmark$ \\
El agua y la humanidad & Rigoberta Menchú & Essay & $\checkmark$ & $\checkmark$ \\
El guardagujas & J.J. Arreola & Kurzgeschichte & -- & $\checkmark$ \\
\hline
\end{tabular}
\end{center}

\vspace{4mm}

%% ABITUR PREP
\secbil{9. Abitur-Vorbereitung}{9. Abitur Preparation}

\vspace{2mm}

\begin{paracol}{2}
\small
\textbf{Vorbereitungsstrategie:}

\vspace{1mm}

\begin{tabular}{|l|l|}
\hline
\rowcolor{lightgray}
\textbf{Phase} & \textbf{Fokus / Zeitraum} \\
\hline
Grundlegend & Alle 4 Themen, Kl. 10--12.1 \\
\hline
Vertiefung & Schwerpunkte lt. Vorabhinw. \\
\hline
Intensiv & Prüfungsformate, Jan--März \\
\hline
Final & Wiederholung, letzte 4 Wo. \\
\hline
\end{tabular}

\vspace{2mm}
\textit{Hinweis:} Die Vorabhinweise benennen jährlich 2 der 4 Themenfelder als Schwerpunkte.

\switchcolumn

\small
\textcolor{english}{\textbf{Preparation Strategy:}}

\vspace{1mm}

\begin{tabular}{|l|l|}
\hline
\rowcolor{lightgray}
\textbf{Phase} & \textbf{Focus / Period} \\
\hline
Foundational & All 4 themes, Gr. 10--12.1 \\
\hline
Deepening & Focus per guidelines \\
\hline
Intensive & Exam formats, Jan--March \\
\hline
Final & Review, last 4 weeks \\
\hline
\end{tabular}

\vspace{2mm}
\textcolor{english}{\textit{Note:} Annual guidelines specify 2 of 4 theme fields as exam focus areas.}

\end{paracol}

\vspace{6mm}

%% ABBREVIATIONS
\vspace{4mm}

\begin{center}
\scriptsize
\begin{tabular}{|l|p{5cm}|p{5cm}|}
\hline
\rowcolor{lightgray}
\textbf{Abk.} & \textbf{Deutsch} & \textbf{English} \\
\hline
GER/GeR & Gemeinsamer Europäischer Referenzrahmen & Common European Framework of Reference (CEFR) \\
\hline
GK & Grundkurs (3 WStd) & Basic Course (3 lessons/week) \\
\hline
LK & Leistungskurs (5 WStd) & Advanced Course (5 lessons/week) \\
\hline
AFB & Anforderungsbereich (I/II/III) & Requirement Level (I/II/III) \\
\hline
WStd & Wochenstunden & Lessons per week \\
\hline
h & Unterrichtsstunden & Teaching hours \\
\hline
\end{tabular}
\end{center}

\vspace{4mm}

\begin{center}
\fcolorbox{spanisch}{lightred}{\parbox{0.95\textwidth}{
\centering
\small
\textbf{Quellenverzeichnis / Sources:}\\[2pt]
\scriptsize
RP Spanisch Qualifikationsphase M-V (2019) | Verknüpfungsbeispiele (2019) | Kommentierte Literaturempfehlungen | Handreichung Sprechprüfung (2020) | Handreichung Paarprüfung Abitur (2025)\\[4pt]
\textit{Hinweis: Wochensequenzen sind eine didaktische Synthese. Der RP gibt keine verbindliche Reihenfolge vor.}\\[4pt]
\textcolor{darkgray}{Lernpfade | Spanisch Oberstufe | Gesamtschule/Gymnasium MV | v1.1 | Januar 2026}
}}
\end{center}

\end{document}
